\section{Sources and current construction boundary}
The baseline syntax reference is ITU-T Recommendation H.263 (02/1998), especially
source format in clause 4.1, start-code byte alignment in clause 4.5,
most-significant-bit transmission in clause 5, picture headers in clause 5.1, group headers in clause 5.2,
macroblock type/presence/motion syntax in clause 5.3, coefficient events and code
tables in clause 5.4.2, motion
compensation in clause 6.1, INTRA DC in clause 5.4.1, inverse quantization and
coefficient clipping in clauses 6.2.1 and 6.2.2, inverse transforms in clause
6.2.4, block reconstruction in clause 6.3 and accuracy in
Annex A. Its optional modes
are outside the selected initial profile. See the
\href{https://www.itu.int/rec/dologin_pub.asp?id=T-REC-H.263-199802-S!!PDF-E&lang=e&type=items}{official recommendation}.
The raw file convention is also described by the
\href{https://ffmpeg.org/ffmpeg-formats.html#rawvideo}{FFmpeg rawvideo documentation}.
MSA instruction semantics follow the
\href{https://s3-eu-west-1.amazonaws.com/downloads-mips/documents/MD00868-1D-MSA64-AFP-01.12.pdf}{MIPS64 SIMD Architecture Module, revision 1.12},
and are checked against the pinned MARS implementation.

The current host interoperability proof uses FFmpeg's baseline H.263 encoder and
an independent decoder invocation, with integer inverse transform selected for
that observation. It recovers eight QCIF pictures from 2194 encoded bytes and
checks their picture types and temporal references. The host encoder is not a
student solution. Its quantized output errors are recorded observations, not the
future acceptance thresholds for a different valid encoder.

This chapter currently contains the raw-input foundation, exact SAD and bounded
full-pixel motion search, baseline interpolation and signed residual constructions.
Six-block prediction now derives target-side origins and stages complete records.
Forward/inverse transform construction and target inverse accuracy checks are now
covered. Quantization now develops a stated forward policy and exact baseline
coefficient reconstruction, including separate INTRA DC. A complete callable block
path now forms transform input, selects levels, reconstructs coefficients and
publishes signed inverse values with clipped pixels. Bounded bit output now
develops equivalent bit and byte-fragment schedules, precise partial-byte
preservation and field-level capacity preflight. Complete coefficient payload
construction now adds scan order, the full compact table, baseline escapes and
private whole-block publication. Complete restricted-baseline macroblock syntax
now combines I/P type, ordinary presence, wrapped motion differences, skip and all
six payloads with private publication. Picture-level construction now stages and publishes both encoded bytes and reconstructed
planes, with explicit row GOB headers and neighbour predictors. A private real-file
driver composes the complete MIPS encoder and independently verifies all eight
generated pictures, quality and finite execution work. The student clip-driver construction now covers bounded short I/O, accepted-reference
publication, exact EOF and descriptor cleanup. Isolated real-file grading includes
independently targeted output failures and decodes every accepted picture. The public
run workflow generates raw samples and independently checks actual target output.
Separate student and instructor distributions now retain the native reading, all twelve
construction stages, public/private feedback and the required full inverse-accuracy
workflow. Their archive and execution receipts record the publication checks.
Six annotated IDE figures now show the worked return states with measured crops,
native numbered explanations and alternative descriptions. Original light-theme
captures, state receipts and color/grayscale proofs are retained. Instructor
implementations and checks remain private.
The authoring/capture plugin stays separate from MARS. Independent accessibility
review and the later substantive editing goal remain open.
